\section{进入与自由进入条件}

本阶段研究 benchmark 的进入模块，并且在固定价格下展开。我们把 Phase 3 已经解出的在位企业值函数视为给定对象，在此基础上推导 entrant values、free-entry conditions、entrant masses 与 entrant-composition vector。全文都把 $(w,p_m,M)$ 视为给定，因此本节还不是 stationary equilibrium existence proof；它只是把后面 market clearing 与 stationary distributions 会用到的 entry objects 先定义清楚并证明其基本性质。

\subsection{Benchmark Fidelity Check}

\paragraph{Step 1：本任务相关的精确 benchmark 对象。}
根据主文稿、benchmark PDF 与项目字典，本阶段相关的 benchmark 对象是：
\begin{itemize}
    \item $V_A(z;M)$、$V_B(z;M)$ 与 $V_C(z;M)$：从 Phase 3 继承而来的在位企业值函数。它们是 dynamic endogenous objects。
    \item $G$：潜在进入者初始能力抽样的分布。这是 entry block 的 primitive distributional object。
    \item $c_{EA},c_{EB},c_{EC}$：以 $A$、$B$、$C$ 形式进入的 entry costs。这些是 primitive cost objects。
    \item $\bar V_E^A(M)$、$\bar V_E^B(M)$、$\bar V_E^C(M)$：分别对应以 $A$、$B$、$C$ 形式进入的 entrant values。这些是 derived entry-value objects。
    \item $n_E^A(M),n_E^B(M),n_E^C(M)$：按组织形式区分的 entrant masses。这些是 endogenous mass objects。
    \item $\mathbf e(M)=\bigl(e_A(M),e_B(M),e_C(M)\bigr)$：entrant-composition vector。这是 derived composition object。
\end{itemize}

\paragraph{Step 2：精确公式、argument list 与经济解释。}
benchmark 文稿把 entrant values 写成
\begin{align*}
\bar V_E^A(M)&=\int_Z V_A(z_0;M)\,G(dz_0)-c_{EA},\\
\bar V_E^B(M)&=\int_Z V_B(z_0;M)\,G(dz_0)-c_{EB},\\
\bar V_E^C(M)&=\int_Z V_C(z_0;M)\,G(dz_0)-c_{EC}.
\end{align*}
benchmark 文稿把 free-entry conditions 写成
\begin{align*}
\bar V_E^A(M) &\le 0, & n_E^A(M)>0 &\implies \bar V_E^A(M)=0,\\
\bar V_E^B(M) &\le 0, & n_E^B(M)>0 &\implies \bar V_E^B(M)=0,\\
\bar V_E^C(M) &\le 0, & n_E^C(M)>0 &\implies \bar V_E^C(M)=0.
\end{align*}
entrant-composition vector 写成
\[
\mathbf e(M)=\bigl(e_A(M),e_B(M),e_C(M)\bigr),
\qquad
e_j(M)=\frac{n_E^j(M)}{n_E^A(M)+n_E^B(M)+n_E^C(M)}.
\]
经济上，$\bar V_E^j(M)$ 表示进入者以组织形式 $j$ 进入时，在对初始能力抽样积分以后得到的期望值函数，并减去相应的 entry cost。$n_E^j(M)$ 表示有多少进入者以形式 $j$ 进入。$\mathbf e(M)$ 则记录进入构成在三种组织形式之间如何分配。

\paragraph{Step 3：本任务能否按 benchmark 原式忠实完成。}
可以，但必须加上一个显式的 domain qualification。benchmark 本身已经足以：
\begin{enumerate}
    \item 由已解出的在位企业值函数与 entrant distribution $G$ 精确地定义 entrant values；
    \item 按 benchmark 原文逐字写出 free-entry conditions；
    \item 在明确写出 entrant masses 非负之后，把 free-entry conditions 改写成等价的 complementarity form；
    \item 在 total entry 严格为正的条件下定义 entrant-composition vector。
\end{enumerate}

\paragraph{Step 4：必须先公开标出的 ambiguity。}
这里有一个必须说明的进入模块歧义。benchmark 的公式
\[
e_j(M)=\frac{n_E^j(M)}{n_E^A(M)+n_E^B(M)+n_E^C(M)}
\]
要求分母严格为正。但 benchmark source 没有单独写出条件
\[
n_E^A(M)+n_E^B(M)+n_E^C(M)>0.
\]
因此在 zero-entry case 下，$\mathbf e(M)$ 并不会自动有定义。本阶段不对这个问题作静默修补，而是只在 total entrant mass 严格为正时定义 entrant-composition vector。

benchmark 文稿里还写了一个 optional extension：
\[
\bar V_E^B(M)=\int_Z V_B(z_0;M)\,G(dz_0)-c_{EB}+\phi_B(M),
\qquad
\phi_B(1)\ge \phi_B(0).
\]
但文稿同时明确说明这个模块 \emph{不属于 benchmark}。因此本阶段不把 $\phi_B(M)$ 纳入正式推导。

\subsection{Part 1：Entrant Values}

\paragraph{Benchmark object restatement.}
对每一个组织形式 $j\in\{A,B,C\}$，benchmark entrant value 定义为：企业以类型 $j$ 开始生命史时的期望在位企业值函数，再减去相应的 entry cost。精确地说，
\begin{align*}
\bar V_E^A(M)&=\int_Z V_A(z_0;M)\,G(dz_0)-c_{EA},\\
\bar V_E^B(M)&=\int_Z V_B(z_0;M)\,G(dz_0)-c_{EB},\\
\bar V_E^C(M)&=\int_Z V_C(z_0;M)\,G(dz_0)-c_{EC}.
\end{align*}
这些对象是 derived dynamic objects。它们不是新的 primitives，也不是新的 Bellman equations。

\paragraph{Economic role.}
在位企业值函数描述的是：企业已经进入行业、并且动态组织选择已经嵌入 Bellman 问题以后，继续经营的 continuation value。进入模块回答的是不同的问题：一个尚未活跃的潜在进入者，如果支付 entry cost 并直接以 $A$、$B$、或 $C$ 形式进入，那么它的期望回报是多少？这个期望回报正是相应的期望在位企业值函数减去相应的进入成本。

\paragraph{Mathematical problem.}
潜在进入者先从分布 $G$ 中抽到初始能力状态 $z_0$，然后比较三种进入形式。对每一个候选进入形式 $j$，entrant value 都是在 $G$ 下对 $V_j(z_0;M)$ 做期望，再减去确定性的成本 $c_{Ej}$。

\paragraph{Admissible set.}
进入时的组织形式选择集合是
\[
\{A,B,C,\mathrm{stay\ out}\}.
\]
本小节只先处理三个 active entry values。outside option 在 free-entry inequalities 里被标准化为 0，因为不进入既不支付 entry cost，也不获得 entrant continuation value。

\paragraph{Assumptions used.}
本节使用以下条件：
\begin{enumerate}
    \item Phase 3 已经证明 $V_A(\cdot;M)$、$V_B(\cdot;M)$ 与 $V_C(\cdot;M)$ 是定义在 $Z$ 上的 bounded continuous functions；
    \item $G$ 是定义在 $Z$ 的 Borel 子集上的 probability measure；
    \item $c_{EA},c_{EB},c_{EC}$ 是有限实数。
\end{enumerate}

\begin{proposition}
在上述条件下，benchmark entrant values $\bar V_E^A(M)$、$\bar V_E^B(M)$ 与 $\bar V_E^C(M)$ 都是良定义且有限的。
\end{proposition}

\begin{proof}
\medskip\noindent\textbf{Step 1：每个 integrand 都是有界且可测的。}
由 Phase 3 可知，每个值函数 $V_j(\cdot;M)$ 都属于 $\mathcal C_b(Z)$。因此每个 $V_j$ 都是定义在 $Z$ 上的 bounded Borel measurable function。

\medskip\noindent\textbf{Step 2：在 $G$ 下的期望是良定义的。}
因为 $G$ 是 probability measure，且 $V_j$ 是 bounded measurable function，所以
\[
\int_Z V_j(z_0;M)\,G(dz_0)
\]
对每个 $j\in\{A,B,C\}$ 都是良定义且有限的。

\medskip\noindent\textbf{Step 3：减去有限的 entry cost。}
由假设，$c_{Ej}$ 是有限的。因此
\[
\bar V_E^j(M)=\int_Z V_j(z_0;M)\,G(dz_0)-c_{Ej}
\]
是有限实数。
\end{proof}

\paragraph{一个有用的界。}
entrant values 从在位企业值函数那里继承了一个简单的有界性结论。对每个 $j\in\{A,B,C\}$，
\begin{align*}
\left|\bar V_E^j(M)\right|
&=
\left|\int_Z V_j(z_0;M)\,G(dz_0)-c_{Ej}\right|\\
&\le
\int_Z |V_j(z_0;M)|\,G(dz_0)+|c_{Ej}|\\
&\le
\|V_j(\cdot;M)\|_\infty+|c_{Ej}|.
\end{align*}
这个界在后面有用，因为它说明：只要 incumbent Bellman block 是有限的，entry block 就不会出现无穷大的对象。

\subsection{Part 2：Free Entry Conditions}

\paragraph{Benchmark object restatement.}
benchmark 文稿把 free-entry conditions 写成
\begin{align*}
\bar V_E^A(M) &\le 0, & n_E^A(M)>0 &\implies \bar V_E^A(M)=0,\\
\bar V_E^B(M) &\le 0, & n_E^B(M)>0 &\implies \bar V_E^B(M)=0,\\
\bar V_E^C(M) &\le 0, & n_E^C(M)>0 &\implies \bar V_E^C(M)=0.
\end{align*}
entrant masses $n_E^A(M),n_E^B(M),n_E^C(M)$ 是 endogenous nonnegative mass objects。

\paragraph{Economic role.}
这些不等式编码了 Hopenhayn 式的 free-entry logic。如果以类型 $j$ 进入的期望净价值严格为正，那么就会有更多质量想按该形式进入，因此 benchmark 的 stationary entry condition 无法成立。所以，只要有正的 entrant mass，以该类型进入的净 entrant value 就必须等于 0。反过来，如果该 entrant value 严格为负，那么 equilibrium 下就不应有正质量按该类型进入。

\paragraph{Mathematical problem.}
这里处理的不是 Bellman step 那种个体层面的 finite-dimensional maximization，而是一个关于 entrant mass 的 equilibrium viability condition。对每一个组织形式，需要共同约束的对象是 entrant value $\bar V_E^j(M)$ 与 entrant mass $n_E^j(M)$。

\paragraph{Admissible set.}
entrant masses 满足
\[
n_E^A(M)\ge 0,\qquad n_E^B(M)\ge 0,\qquad n_E^C(M)\ge 0.
\]
质量非负在经济上本来就是 mass 的定义内容。但这里要把它显式写出来，因为把 benchmark 的 implication form 翻译成 complementarity form 时必须用到这一点。

\begin{proposition}
如果 entrant masses 非负，那么对每个 $j\in\{A,B,C\}$，benchmark 的 free-entry condition 等价于
\[
\bar V_E^j(M)\le 0,\qquad n_E^j(M)\ge 0,\qquad n_E^j(M)\bar V_E^j(M)=0.
\]
\end{proposition}

\begin{proof}
固定一个组织形式 $j$。

\medskip\noindent\textbf{Step 1：benchmark 的 implication form 蕴含 complementarity.}
假设 benchmark 条件
\[
\bar V_E^j(M)\le 0
\]
以及
\[
n_E^j(M)>0\implies \bar V_E^j(M)=0.
\]
同时再加入 $n_E^j(M)\ge 0$，因为 entrant mass 不可能为负。分两种情况。

如果 $n_E^j(M)=0$，那么自动有
\[
n_E^j(M)\bar V_E^j(M)=0.
\]
如果 $n_E^j(M)>0$，那么由 benchmark implication 可得
\[
\bar V_E^j(M)=0,
\]
因此同样有
\[
n_E^j(M)\bar V_E^j(M)=0.
\]
所以 complementarity 的乘积条件成立。

\medskip\noindent\textbf{Step 2：complementarity 蕴含 benchmark 的 implication form.}
现在反过来假设
\[
\bar V_E^j(M)\le 0,\qquad n_E^j(M)\ge 0,\qquad n_E^j(M)\bar V_E^j(M)=0.
\]
如果 $n_E^j(M)>0$，那么把等式
\[
n_E^j(M)\bar V_E^j(M)=0
\]
除以正数 $n_E^j(M)$，得到
\[
\bar V_E^j(M)=0.
\]
这正是 benchmark 的 implication：
\[
n_E^j(M)>0\implies \bar V_E^j(M)=0.
\]
因此，一旦把 entrant mass 的非负性显式写出来，这两种写法就是等价的。
\end{proof}

\paragraph{complementarity 的经济含义是什么？}
complementarity form 表示：同一个组织进入边际上，不能同时出现“正的 entrant mass”和“严格负的 entrant value”。如果质量为正，那么 entrant value 必须为零；如果 entrant value 为负，那么 entrant mass 必须为零。这正是本 benchmark 里 free entry 的全部经济内容。

\subsection{Part 3：Entrant Masses 与 Entrant-Composition Vector}

\paragraph{Benchmark object restatement.}
benchmark 文稿先定义按类型区分的 entrant masses，
\[
n_E^A(M),\qquad n_E^B(M),\qquad n_E^C(M),
\]
然后定义 entrant-composition vector
\[
\mathbf e(M)=\bigl(e_A(M),e_B(M),e_C(M)\bigr),
\qquad
e_j(M)=\frac{n_E^j(M)}{n_E^A(M)+n_E^B(M)+n_E^C(M)}.
\]

\paragraph{Economic role.}
$n_E^j(M)$ 描述的是不同组织形式上的进入规模。$\mathbf e(M)$ 描述的是进入构成如何在不同类型之间分配。这与 incumbent reorganization 在概念上不同。incumbent reorganization 告诉我们已有企业如何在组织形式之间切换或退出；entry composition 告诉我们新进入行业的企业第一次被观察到时，以哪种组织形式进入。

\paragraph{Mathematical problem.}
我们要回答的是：$\mathbf e(M)$ 在什么条件下有定义，并且它有什么性质。答案完全取决于分母
\[
N_E(M):=n_E^A(M)+n_E^B(M)+n_E^C(M)
\]
的符号。这里的记号 $N_E(M)$ 只是为了方便证明而引入的 auxiliary notation；它只是 total entrant mass。

\paragraph{Admissible set.}
因为 entrant masses 非负，
\[
N_E(M)\ge 0.
\]
benchmark 比值定义要成立，需要更强的条件
\[
N_E(M)>0.
\]
如果 $N_E(M)=0$，那么 benchmark source 并没有定义 $\mathbf e(M)$。

\begin{proposition}
如果
\[
N_E(M)=n_E^A(M)+n_E^B(M)+n_E^C(M)>0,
\]
那么 entrant-composition vector 是良定义的，并满足
\[
0\le e_j(M)\le 1
\qquad\text{对每个 }j\in\{A,B,C\},
\]
以及
\[
e_A(M)+e_B(M)+e_C(M)=1.
\]
\end{proposition}

\begin{proof}
\medskip\noindent\textbf{Step 1：这些比值是良定义的。}
如果 $N_E(M)>0$，那么
\[
e_j(M)=\frac{n_E^j(M)}{N_E(M)}
\]
里的分母严格为正，所以每个比值都有定义。

\medskip\noindent\textbf{Step 2：每个分量都落在 $[0,1]$ 内。}
因为 entrant masses 非负，
\[
0\le n_E^j(M)\le N_E(M).
\]
把不等式两边都除以正数 $N_E(M)$，得到
\[
0\le e_j(M)\le 1.
\]

\medskip\noindent\textbf{Step 3：各分量之和等于 1。}
直接代入得
\begin{align*}
e_A(M)+e_B(M)+e_C(M)
&=
\frac{n_E^A(M)}{N_E(M)}
+
\frac{n_E^B(M)}{N_E(M)}
+
\frac{n_E^C(M)}{N_E(M)}\\
&=
\frac{n_E^A(M)+n_E^B(M)+n_E^C(M)}{N_E(M)}\\
&=
\frac{N_E(M)}{N_E(M)}\\
&=1.
\end{align*}
\end{proof}

\paragraph{为什么 entrant composition 在概念上不同于 incumbent reorganization？}
benchmark 文稿对这一点说得很明确，而且这个区分在理论与经验上都重要。incumbent transition matrix $\mathbf P(M)$ 描述的是既有企业做了什么：留在原组织形式、切换到其他组织形式、还是退出。entrant-composition vector $\mathbf e(M)$ 描述的是另一种对象：那些在期初并不是 active incumbents、而是在本期从行业外部进入的企业，其中有多少比例以 $A$ 开始，有多少以 $B$ 开始，有多少以 $C$ 开始。这是不同的 margin，不能混为一谈。type-$B$ entry 上升，并不自动意味着 incumbent 向 type $B$ 的 switching 也更多；反过来也是一样。

\subsection{What did we just do, and why does it matter?}

我们把 Phase 3 的 incumbent value functions 转换成了 entrant values。这个步骤在形式上很简单，但在经济上很关键。它说明进入回报并不是来自一个单独的 entrant production problem；相反，它来自潜在进入者直接进入某个 benchmark 组织形式后所获得的期望 continuation value，再减去相应的 entry cost。

接着我们加入 free-entry logic。这一步重要，因为仅有 incumbent values 还不足以 pin down entry。模型还必须对每个组织形式上究竟能有多少 entrant mass 获利进入给出 equilibrium restriction。free-entry inequalities 与 complementarity conditions 正是把 entrant profitability 与 entrant masses 连起来的环节。

entrant-composition vector 之所以重要，是因为它和 incumbent transition matrix 不是同一个对象。transition matrix 描述的是已经在市场中的企业如何重新组织。entrant-composition vector 描述的是从市场外进入的新企业，其组织形式的构成。这个区分在数据里必不可少，在理论里也同样必不可少。如果把两者混在一起，模型就无法区分 MAH 究竟主要是通过 incumbent reallocation 改变组织边界，还是主要通过 shifting new entry 改变组织边界，还是两者同时发生。

这里还有一个我们明确写出的 domain issue，而不是把它藏起来。entrant-composition vector 只有在 total entrant mass 严格为正时才有定义。benchmark source 写出了公式，但没有单独拼出这个 positivity condition。把这个条件显式写出来，可以保证后续 stationary-distribution block 与 equilibrium block 的数学写法保持清楚。

\subsection{End-of-Section Recap}

\paragraph{What has been established mathematically.}
在固定 $(w,p_m,M)$ 下，benchmark entrant values 是已解出的 incumbent value functions 在 entrant distribution $G$ 下的有限期望，并减去 entry costs。benchmark free-entry conditions 可以按 source 原式写出；一旦把 entrant masses 的非负性显式写出，它们也可以等价地写成 complementarity form。如果 total entrant mass 严格为正，那么 entrant-composition vector $\mathbf e(M)$ 良定义、每个分量都落在 $[0,1]$ 内，并且三者之和为 1。

\paragraph{What assumptions were used.}
这里使用了 Phase 3 的结果：$V_A$、$V_B$ 与 $V_C$ 是 bounded continuous functions；还使用了 benchmark entrant distribution $G$、有限的 entry costs $c_{EA},c_{EB},c_{EC}$、以及 entrant masses 非负。为了定义 $\mathbf e(M)$，我们额外要求 total entrant mass 严格为正。

\paragraph{What remains to be derived next.}
下一阶段需要把 incumbent objects 与 entry objects 一起带入 market clearing，先处理 qualified contract-manufacturing market，再处理 labor market。之后的 stationary-distribution phase 还需要把 entrant masses 与 optimal policy behavior 结合起来，定义 invariant measures。
